% Tutorial set: 02
% Source: T2.pdf, Tutorial 2 (Image Enhancement in the Spatial Domain), Questions 1--3
% Reference: T2_Reference solutions.pptx, used for answer checking and qualitative plots
% Session date: 2026-09-01
% Human verified: Yes

\clearpage
\subsection{Tutorial 02：Spatial-domain Image Enhancement}
\label{tutorial:SC4061-TUT02}

\exercisemeta
  {SC4061-TUT02}
  {2026-09-01}
  {T2.pdf，Questions 1--3；T2\_Reference solutions.pptx 仅用于核对}
  {三题均已作答、讨论并完成订正}
  {已确认（2026-09-01）}

\exercisequestion{SC4061-TUT02-Q01}{Monitor Gamma Calibration}

\textbf{对应讲义：}
\relatedtopic{SC4061-L03-T02}{Contrast、Gamma 与 Histogram Equalisation}

\subsubsection*{题目}

Monitor（显示器）的输入 gray level（灰度值）$r$ 与 photosensitive device（光敏设备）测得的输出 $s$ 被假设满足 power law（幂律）
\[
  s=cr^{\gamma}.
\]
实验给出 $100\mapsto0.24$ 与 $200\mapsto0.78$。（a）求 $\gamma$ 并判断曲线形状；（b）若输出单位改为 $24$ 与 $78$，判断估计的 $\gamma$ 是否改变；（c）说明该校准模型的不足。

\begin{checkedsolution}
两组方程相除可消去未知比例系数 $c$：
\[
  \frac{s_2}{s_1}=\left(\frac{r_2}{r_1}\right)^{\gamma},
  \qquad
  \gamma=\frac{\ln(s_2/s_1)}{\ln(r_2/r_1)}
  =\frac{\ln(13/4)}{\ln2}
  \approx\boxed{1.7004}.
\]
$\gamma>1$，故曲线 increasing and convex（单调递增且向上凸）。若所有输出统一乘以常数 $k$，ratio $ks_2/(ks_1)$ 不变，因此 $\gamma$ 不变，只会改变 $c$。

真实显示器通常不严格服从单一 power law；尤其在 $r=0$ 时仍可能存在 black offset（黑电平偏置）。较合理的简单模型可写为
\[
  s=b+cr^{\gamma},
\]
或使用 piecewise model / LUT（分段模型/查找表）。Ambient light、sensor calibration、screen position 与 colour channel 也会影响测量。
\end{checkedsolution}

\textbf{检查点：}改变统一的输入或输出单位不会改变由 ratios 求得的 $\gamma$；black offset 改变的是模型结构，不能仅靠重新估计 $c$ 修复。

\exercisequestion{SC4061-TUT02-Q02}{Step Edge 的 Box 与 Gaussian Smoothing}

\textbf{对应讲义：}
\relatedtopic{SC4061-L03-T04}{Box、Gaussian 与 Median Filtering}

\subsubsection*{题目}

一幅 grayscale image（灰度图像）左半为黑色、右半为白色，其横向 profile（剖面）是从 $0$ 跳至 $255$ 的 step edge（阶跃边缘），并假设 spatial continuous domain（连续空间域）。（a）画原图 histogram；（b）经 $L\times L$ box filter 后画输出 profile 与 histogram；（c）改用 Gaussian smoothing 后重复作图；（d）给出 MATLAB verification（验证）代码。

\begin{checkedsolution}
设edge位于 $x=0$：
\[
  r(x)=
  \begin{cases}
    0,&x<0,\\
    255,&x\ge0.
  \end{cases}
\]
原图只有两个 gray levels，normalized histogram 为
\[
  p_r(t)=\tfrac12\delta(t)+\tfrac12\delta(t-255).
\]

图像沿 $y$ 方向不变，因此 $L\times L$ box filter 可按一维宽度 $L$ 分析：
\[
  g_B(x)=
  \begin{cases}
    0,&x\le-L/2,\\[1mm]
    255\left(x/L+1/2\right),&-L/2<x<L/2,\\[1mm]
    255,&x\ge L/2.
  \end{cases}
\]
Step 被展开成宽度为 $L$ 的 linear ramp（线性斜坡）。大片纯黑/纯白区域仍产生 $0$ 与 $255$ 的 peaks；ramp slope 为常数 $255/L$，故中间 gray levels 对应相同空间宽度，在 histogram 中形成 uniform plateau（均匀平台）。

Gaussian kernel 与step的卷积为 Gaussian CDF：
\[
  g_G(x)=255\Phi\!\left(\frac{x}{\sigma}\right)
  =\frac{255}{2}\left[1+\operatorname{erf}\!\left(\frac{x}{\sqrt2\sigma}\right)\right].
\]
其profile为smooth S-curve（平滑S形曲线）。由变量变换
\[
  p_g(t)\propto\left|\frac{dx}{dt}\right|
  =\frac{1}{|dg_G/dx|},
\]
profile在中间斜率最大，因此中间gray levels较少；两端斜率小，因此histogram两端较高。Ideal Gaussian 具有无限support，只会渐近接近 $0$ 与 $255$；实际截断kernel会保留远离edge的纯黑/纯白区域，并在两端形成明显peaks。

\begin{figure}[htbp]
  \centering
  \resizebox{0.98\textwidth}{!}{\begin{tikzpicture}[x=0.9cm,y=0.82cm,font=\scriptsize]
  \tikzset{
    axis/.style={->,semithick,gray!75!black},
    signal/.style={very thick,noteBlue},
    hist/.style={very thick,orange!85!black}
  }

  % Profiles
  \begin{scope}[shift={(0,3.7)}]
    \node[font=\small\bfseries] at (1.8,2.1) {Original step};
    \draw[axis] (0,0) -- (3.7,0) node[right] {$x$};
    \draw[axis] (0,0) -- (0,1.8) node[above] {$r$};
    \draw[signal] (0,0) -- (1.8,0) -- (1.8,1.5) -- (3.5,1.5);
    \node[left] at (0,0) {$0$};
    \node[left] at (0,1.5) {$255$};
  \end{scope}

  \begin{scope}[shift={(5.2,3.7)}]
    \node[font=\small\bfseries] at (1.8,2.1) {Box-filtered profile};
    \draw[axis] (0,0) -- (3.7,0) node[right] {$x$};
    \draw[axis] (0,0) -- (0,1.8) node[above] {$g_B$};
    \draw[signal] (0,0) -- (1.15,0) -- (2.45,1.5) -- (3.5,1.5);
    \draw[<->,gray!75!black] (1.15,-0.28) -- node[below] {$L$} (2.45,-0.28);
  \end{scope}

  \begin{scope}[shift={(10.4,3.7)}]
    \node[font=\small\bfseries] at (1.8,2.1) {Gaussian-filtered profile};
    \draw[axis] (0,0) -- (3.7,0) node[right] {$x$};
    \draw[axis] (0,0) -- (0,1.8) node[above] {$g_G$};
    \draw[signal,smooth,samples=80,domain=0.15:3.45]
      plot (\x,{1.5/(1+exp(-4*(\x-1.8)))});
    \draw[dashed,gray!60] (1.8,0) -- (1.8,0.75);
    \node[below] at (1.8,0) {$x_0$};
  \end{scope}

  % Histograms
  \begin{scope}[shift={(0,0)}]
    \node[font=\small\bfseries] at (1.8,2.1) {Original histogram};
    \draw[axis] (0,0) -- (3.7,0) node[right] {gray level};
    \draw[axis] (0,0) -- (0,1.8) node[above] {$h$};
    \draw[hist,->] (0.55,0) -- (0.55,1.45);
    \draw[hist,->] (3.05,0) -- (3.05,1.45);
    \node[below] at (0.55,0) {$0$};
    \node[below] at (3.05,0) {$255$};
  \end{scope}

  \begin{scope}[shift={(5.2,0)}]
    \node[font=\small\bfseries] at (1.8,2.1) {Box-filtered histogram};
    \draw[axis] (0,0) -- (3.7,0) node[right] {gray level};
    \draw[axis] (0,0) -- (0,1.8) node[above] {$h$};
    \draw[hist,->] (0.55,0) -- (0.55,1.45);
    \draw[hist,->] (3.05,0) -- (3.05,1.45);
    \draw[hist] (0.55,0.38) -- (3.05,0.38);
    \node[below] at (0.55,0) {$0$};
    \node[below] at (3.05,0) {$255$};
    \node[above] at (1.8,0.38) {uniform intermediate levels};
  \end{scope}

  \begin{scope}[shift={(10.4,0)}]
    \node[font=\small\bfseries] at (1.8,2.1) {Gaussian-filtered histogram};
    \draw[axis] (0,0) -- (3.7,0) node[right] {gray level};
    \draw[axis] (0,0) -- (0,1.8) node[above] {$h$};
    \draw[hist,->] (0.55,0) -- (0.55,1.45);
    \draw[hist,->] (3.05,0) -- (3.05,1.45);
    \draw[hist,smooth,samples=60,domain=0.55:3.05]
      plot (\x,{0.28+0.55*(\x-1.8)^2});
    \node[below] at (0.55,0) {$0$};
    \node[below] at (3.05,0) {$255$};
    \node[above] at (1.8,0.28) {few mid-gray pixels};
  \end{scope}
\end{tikzpicture}
}
  \caption{Q2 的profile与histogram对照。依据\textit{T2} Question 2及\textit{T2 Reference Solutions} slides 7--8重绘；histogram纵轴只表示定性频数。}
  \label{fig:sc4061-tut02-filter-profiles}
\end{figure}

\textbf{MATLAB verification：}
\begin{verbatim}
N = 512; L = 31; sigma = 5;
I = zeros(N); I(:, N/2+1:end) = 255;
B = imfilter(I, ones(L,L)/(L^2), 'replicate', 'same');
G = imgaussfilt(I, sigma, 'FilterSize', L, ...
                'Padding', 'replicate');

x = 1:N; row = N/2;
subplot(2,2,1); plot(x, I(row,:), x, B(row,:));
subplot(2,2,2); histogram(B(:), 0:256);
subplot(2,2,3); plot(x, I(row,:), x, G(row,:));
subplot(2,2,4); histogram(G(:), 0:256);
\end{verbatim}
\end{checkedsolution}

\textbf{检查点：}Profile 的纵轴是输出intensity，histogram的纵轴是pixel frequency；$g_B(0)=127.5$ 不等于“histogram在$0$处变为$127.5$”。

\exercisequestion{SC4061-TUT02-Q03}{Median Filter 对 Spikes 与 Step Edge 的作用}

\textbf{对应讲义：}
\relatedtopic{SC4061-L03-T04}{Box、Gaussian 与 Median Filtering}

\subsubsection*{题目}

给定1D signal
\[
  [2,6,53,10,7,4,107,103,104,109,57,102].
\]
其中包含两个spikes（尖峰）和一个step edge。（a）画出signal；（b）采用endpoint replication（端点复制）边界，对其施加3-element median filter并写出输出；（c）对输出再次施加同一filter；（d）说明若改用Gaussian filtering，spikes与edge如何变化；（e）说明何时median filtering更合适。

\begin{checkedsolution}
三元素中值滤波为
\[
  y_i=\operatorname{median}(x_{i-1},x_i,x_{i+1}),
  \qquad x_0=x_1,\quad x_{13}=x_{12}.
\]
第一次输出为
\[
  \boxed{[2,6,10,10,7,7,103,104,104,104,102,102]}.
\]
对该结果再次滤波，输出不再改变。两个isolated spikes被移除，而low-to-high step仍被保留。

Gaussian filtering使用weighted average（加权平均）：spike amplitude会减小，但异常值也会扩散到邻近位置；step edge会变宽、斜率降低。Median filtering更适合salt-and-pepper noise（椒盐噪声）或少量extreme outliers，并通常更能保留sharp edges。

Median能恢复正常值的条件是窗口内被污染样本少于一半；若3个元素中至少2个异常，median也会失败。此外，“第二次不再变化”只对本题signal成立，并非median filter普遍只需执行一次。
\end{checkedsolution}

\textbf{检查点：}Median filter是nonlinear filter（非线性滤波器），不能用convolution theorem分析；其robustness来自order statistic（次序统计量），而不是对异常值赋予较小线性权重。

\subsubsection*{本次 Tutorial 收口}

Q1训练power-law calibration与模型失配；Q2用profile slope解释smoothing后的histogram；Q3比较median与Gaussian对outlier和edge的不同影响。复习时优先掌握：ratio消去尺度、histogram density与profile slope成反比，以及median成功需要窗口内正常样本占多数。
